\textbf{T1: Jumper (10 pts)}

Two uniform rods, each of mass $m$ and length $\ell$, are connected at one end
by a hinge equipped with a torsional spring. The spring exerts equal and
opposite restoring torques $\tau$ on the rods about the hinge,
\[
\tau = 2k\theta,
\]
where $2\theta$ is the angle between the rods (in radians), and $k$
($k \gg mg\ell$) is the torsional spring constant.

\begin{center}\includegraphics[width=0.37\linewidth]{figures/EuPhO_2026_Q1_fig1.png}\end{center}

The system is placed on a flat horizontal surface with the free ends of the
rods resting on the floor and the hinge above them. The hinge is then pushed
downward until the rods lie flat on the floor in opposite directions,
corresponding to $\theta = \pi/2$. Friction between the rods and the floor is
negligible.

Jumper jumps after the hinge is released. Determine the maximum height $h$
reached by the center of mass of the system to a precision better than 1\%.
